\scanpage{224}
\mspage{174}
% L1–2
celle justement qui associe à tout $x\in X$ le point marqué $1_x$ de $\Omega(X;x,x)$ (``lacet vide en $x$'').

% L3
La question est donc celle-ci :

% L4–6
On a une $W$-fibration $\Omega\longrightarrow X$, avec une section donnée \inserted{$\sigma$}. Cela définit-il \deleted{une application} foncteur canonique
% L7
\[
X\longrightarrow\mathrm{Hot}_{\bullet W},
\]
% L8
transformant toute flèche en iso ? 

% L9
\deleted{\unsure{L'i}} L'idée naturelle est de prendre
% L10
\[
x\longmapsto\mathrm{hot}_{\bullet W}(\deleted{X}\Omega_{/x})
\]
% L11–14
en ponctuant $\deleted{X}\Omega_{/x}$ par le pt marqué $\sigma(x)$ de $\deleted{X}\Omega_x\hookrightarrow\Omega_{/x}$. L'ennui, c'est que les morphismes de transition $\deleted{X}\Omega_{/x}\unsure{-|\to}\deleted{X}\Omega_{/x'}$ (pour $x\xrightarrow{u}x'$ dans $X$) ne sont pas compatibles avec les pts marqués.

% L15–19
\deleted{Mais. \deleted{\unsure{si on a une flèche} \deleted{\unsure{pour}} \unsure{savoir $\sigma(x)$) elle est}} Si on regarde $\sigma(X_{/x'})$, c'est une sous-catégorie $W$-asphérique de $\Omega_{/x'}$ qui contient et le pt marqué de $\Omega_{/x'}$, et l'image du point marqué de $\Omega_{/x}$. Cela ne permet-il pas de définir la flèche cherchée ?}

% L20–21
\deleted{Une autre façon de \unsure{le|se} présenter est celle-ci. On regarde, pour $x$ fixé, dans $X$}
% L22–24 left
\deleted{\[
\unsure{\sigma}(X_{/x})\longrightarrow\Omega_{/x}
\]}
% L22–24
\deleted{\diagram{\begin{tikzcd}[ampersand replacement=\&,column sep=large]
X_{/x}\arrow[r,"\sigma"]\arrow[d,"{\in W}"]\arrow[dr,squiggly,dashed] \& \Omega_{/x}\arrow[d,dotted]\\
e\arrow[r,dashed,no head] \& \left(\begin{gathered}\qquad\nearrow\,\Omega_{/x}\\[-3pt]X_{/x}\qquad\\[-3pt]\qquad\searrow\,e\end{gathered}\right)\overset{\mathrm{def}}{=}\widetilde{\Omega}(x)
\end{tikzcd}}}
% archive sheet number: 224
